\documentclass[aspectratio=169,12pt]{beamer}
\usetheme{Madrid}
\usecolortheme{dolphin}
\usefonttheme{professionalfonts}
\setbeamertemplate{navigation symbols}{}
\setbeamertemplate{footline}{}

\usepackage[T1]{fontenc}
\usepackage[utf8]{inputenc}
\usepackage{graphicx}
\usepackage{amsmath,amssymb}
\usepackage{hyperref}
\usepackage{siunitx}
\usepackage{physics}
\usepackage{braket}

\graphicspath{{../../figures/}}

\title[Module B4]{Module B4: Distribution quantique de clés}
\author{JAMOTTE Maxime, SCHOONEN Cédric}
\institute{Digital Learning Hub}
\date{}

\begin{document}

\begin{frame}
  \titlepage
\end{frame}

\begin{frame}{Distribution quantique de clés (voir JN)}
  \begin{itemize}
    \setlength{\itemsep}{0.8em}
    \item Alice prépare des qubits aléatoirement dans la base $Z$ ou $X$.\\~\\
    \item Bob mesure aléatoirement dans la base $Z$ ou $X$.\\~\\
    \item Ils partagent leurs choix de bases et ne gardent que les résultats où les bases coïncident.\\~\\
    \item Résultat: une clé secrète partagée.\\~\\
    \item Toute interception (Eve) introduit des erreurs détectables.\\~\\
    \item Voir le Jupyter Notebook pour l'implémentation complète avec Qiskit.
  \end{itemize}
\end{frame}

\begin{frame}{Fin du module B4}
  \centering
  \vfill
  Module optionnel recommandé.
  \vfill
\end{frame}

\end{document}
